$\beta$ dependent eloss cut ($E_{\text{loss}}^{\text{threshold}} = p_0 \beta^{p_1}$) was first introduced in \href{https://phenix-intra.sdcc.bnl.gov/phenix/WWW/p/draft/chujo/PPG026/AnalysisNote/}{AN187} to remove the background in TOFe. It was also employed in this analysis with parameters $p_0 = 0.0005, p_1 = -2.6$ for real data and $p_0 = 0.0016, p_1 = -2.6$ for MC for all datasets. Examples of $\beta$ vs $E_{\text{loss}}$ distributions and $E_{\text{loss}}$ cut threshold (red line) are shown on Fig. \ref{fig:eloss_cut}. Data with $E_{\text{loss}}$ lower than $E_{\text{loss}}^{\text{threshold}}$ was removed from this analysis.


\begin{figure}[h!]
	\centering
   \begin{subfigure}{0.45\textwidth}
      \includegraphics[width=\textwidth]{ElossFit/Run14HeAu200.pdf}
      \subcaption{}
   \end{subfigure}
   \hfill
   \begin{subfigure}{0.45\textwidth}
      \includegraphics[width=\textwidth]{ElossFit/SimRun14HeAu200.pdf}
      \subcaption{}
   \end{subfigure}
   \caption[Examples of $\beta$ vs $E_{\text{loss}}$ distributions and $E_{\text{loss}}$ cut]
   {Examples of $\beta$ vs $E_{\text{loss}}$ distributions and $E_{\text{loss}}$ cut (red line) in Run14 \HeAu for real data (a) and for MC (b)}
   \label{fig:eloss_cut}
\end{figure}
